\documentclass[a4paper]{article}

\usepackage[spanish]{babel}
\usepackage{graphicx}
\usepackage{amsmath, amssymb}
\usepackage[margin=2cm]{geometry}
\usepackage{fancyhdr}
\usepackage{enumerate}
\usepackage[shortlabels]{enumitem}
\usepackage{parskip}
\usepackage[most]{tcolorbox}
\usepackage[hidelinks]{hyperref}
\usepackage{float}
\usepackage{booktabs}
\usepackage{tikz}
\usetikzlibrary{arrows.meta, angles, quotes, decorations.pathmorphing, babel}

% cabecera
\pagestyle{fancy}
\fancyhead[l]{Física Matemática I}
\fancyhead[c]{Notas de Clase 20260612}
\fancyhead[r]{\today}
\fancyfoot[c]{\thepage}
\setlength{\headheight}{14.5pt}
\renewcommand{\headrulewidth}{0.2pt}

\begin{document}
\pagenumbering{gobble}

\begin{titlepage}
  \thispagestyle{empty}
  \centering
  {\Large Universidad de Antioquia\par}
  \vspace{2cm}
  {\huge\bfseries Física Matemática I\par}
  \vspace{0.4cm}
  {\Large Clase 20260612\par}
  \vspace{0.8cm}
  {\large Armónicos esféricos, función de Green y ecuación de Laguerre\par}
  \vspace{1.7cm}
  {\large Autor:\par}
  \vspace{0.2cm}
  {\large Daniel Soto Villada\par}
  {\large Estudiante de Astronomía\par}
  \vfill
  {\large \today\par}
\end{titlepage}

\pagenumbering{roman}
\setcounter{page}{1}
\newpage
\tableofcontents
\newpage
\pagenumbering{arabic}

\section*{Armónicos esféricos, función de Green y ecuación de Laguerre}
\addcontentsline{toc}{section}{Armónicos esféricos, función de Green y ecuación de Laguerre}


\textit{Transcrito y organizado a partir de las fotografías del tablero (clase del 12/06/2026), con explicaciones adicionales para dar contexto y completar los pasos intermedios.}

\medskip


\subsection{Introducción}


En esta clase se resuelven dos problemas de electrostática con simetría esférica mediante \textbf{separación de variables en armónicos esféricos}. Para el segundo problema se necesita además el \textbf{teorema de adición de armónicos esféricos} y el método de la \textbf{función de Green} para la ecuación de Poisson, que conduce a la \textbf{expansión multipolar} en polinomios de Legendre. Al final se introduce la \textbf{ecuación de Laguerre}, cuyas soluciones ---junto con los armónicos esféricos--- son la base de las funciones de onda del átomo de hidrógeno (de ahí su lugar en este curso).

\medskip


\section{Repaso: ecuación de Laplace en coordenadas esféricas}


La solución general de $\nabla^2\phi=0$ en coordenadas esféricas $(r,\theta,\varphi)$ es:

\[
\phi(r,\theta,\varphi)=\sum_{\ell=0}^{\infty}\sum_{m=-\ell}^{\ell}\left[A_{\ell m}\,r^{\ell}+\frac{B_{\ell m}}{r^{\,\ell+1}}\right]Y_\ell^m(\theta,\varphi)
\]


donde los \textbf{armónicos esféricos} son

\[
Y_\ell^m(\theta,\varphi)=\sqrt{\frac{2\ell+1}{4\pi}\,\frac{(\ell-m)!}{(\ell+m)!}}\;P_\ell^{m}(\cos\theta)\,e^{im\varphi}
\]


Los dos casos que se usan en el Ejemplo de la sección 2 (ambos con $m=0$, sin dependencia en $\varphi$):

\[
Y_0^0=\sqrt{\frac{1}{4\pi}}=\text{cte}, \qquad\qquad Y_1^0=\sqrt{\frac{3}{4\pi}}\,P_1(\cos\theta)=\sqrt{\frac{3}{4\pi}}\cos\theta
\]


\textbf{Condición de regularidad en el origen:} si la región de interés incluye $r=0$ (problema \textbf{interior}), se debe tener $B_{\ell m}=0$ para todo $\ell,m$ (de lo contrario $\phi\to\infty$ cuando $r\to0$). Por lo tanto:

\[
\phi_{\text{int}}(r,\theta,\varphi)=\sum_{\ell,m}A_{\ell m}\,r^{\ell}\,Y_\ell^m(\theta,\varphi)\tag{1}
\]


\medskip


\section{\texorpdfstring{Ejemplo resuelto: cascarón esférico con $\phi(R,\theta)=V_0\sin^2(\theta/2)$}{Ejemplo resuelto: cascarón esférico con potencial de frontera}}


\textbf{Enunciado:} determinar el campo eléctrico y el potencial en el \textbf{interior} de un cascarón esférico de radio $R$, dado que el potencial sobre su superficie es

\[
\phi(R,\theta)=V_0\sin^2\!\left(\frac{\theta}{2}\right)
\]


\begin{center}
\begin{tikzpicture}[scale=1.05, >=Stealth]
  \coordinate (O) at (0,0);
  \coordinate (P) at (42:2.2);
  \coordinate (Z) at (0,2.2);
  \draw[fill=blue!5, draw=blue!55!black, thick] (O) circle (2.2);
  \draw[->, thick] (O) -- (0,2.9) node[above] {$z$};
  \draw[->, thick] (O) -- (3.0,0) node[right] {$x$};
  \draw[->, very thick, red!70!black] (O) -- (P) node[midway, above left] {$R$};
  \draw[dashed] (P) -- (0,2.2);
  \pic[draw=black, ->, "$\theta$", angle radius=0.8cm, angle eccentricity=1.25] {angle=P--O--Z};
  \fill (O) circle (1.3pt) node[below left] {$O$};
  \fill (P) circle (1.3pt) node[above right] {$\phi(R,\theta)=V_0\sin^2(\theta/2)$};
  \node[align=center, blue!55!black] at (0,-2.65) {cascarón esférico\\superficie $r=R$};
\end{tikzpicture}
\end{center}


\subsection{Planteamiento}


Como la región de interés (interior) incluye el origen, se usa directamente (1):

\[
\phi(r,\theta)=\sum_{\ell,m}A_{\ell m}\,r^{\ell}\,Y_\ell^m(\theta,\varphi)
\]


\subsection{Expansión del potencial de frontera en armónicos esféricos}


En $r=R$ debe cumplirse:

\[
V_0\sin^2\!\left(\frac{\theta}{2}\right)=\sum_{\ell,m}A_{\ell m}\,R^{\ell}\,Y_\ell^m(\theta,\varphi)
\]


Usando la identidad trigonométrica $\sin^2(\theta/2)=\dfrac{1-\cos\theta}{2}$:

\[
V_0\sin^2\!\left(\frac{\theta}{2}\right)=\frac{V_0}{2}-\frac{V_0}{2}\cos\theta
\]


Esta expresión no depende de $\varphi$ y solo contiene un término constante y un término en $\cos\theta$, es decir, únicamente los armónicos $Y_0^0$ y $Y_1^0$ intervienen. Invirtiendo las relaciones de la sección 1:

\[
1=\sqrt{4\pi}\;Y_0^0\,,\qquad\qquad \cos\theta=\sqrt{\frac{4\pi}{3}}\;Y_1^0
\]


Sustituyendo:

\[
V_0\sin^2\!\left(\frac{\theta}{2}\right)=\frac{V_0\sqrt{4\pi}}{2}\left[Y_0^0-\frac{1}{\sqrt3}\,Y_1^0\right]
\]


\subsection{\texorpdfstring{Cálculo de los coeficientes $A_{\ell m}$}{Cálculo de los coeficientes A lm}}


Se multiplica ambos lados por $(Y_{\ell'}^{m'})^*(\theta,\varphi)$ y se integra sobre el ángulo sólido, usando la ortonormalidad de los armónicos esféricos:

\[
\int Y_\ell^m\,(Y_{\ell'}^{m'})^*\,d\Omega=\delta_{\ell\ell'}\,\delta_{mm'}
\]


Esto da $A_{\ell'm'}R^{\ell'}=\dfrac{V_0\sqrt{4\pi}}{2}\left[\delta_{\ell'0}\delta_{m'0}-\dfrac{1}{\sqrt3}\delta_{\ell'1}\delta_{m'0}\right]$, es decir:

\[
\boxed{A_{\ell m}=\frac{V_0\sqrt{4\pi}}{2}\,R^{-\ell}\left[\delta_{\ell0}\delta_{m0}-\frac{1}{\sqrt3}\,\delta_{\ell1}\delta_{m0}\right]}\tag{2}
\]


Solo dos coeficientes son distintos de cero:

\[
A_{00}=V_0\sqrt{\pi}\,,\qquad\qquad A_{10}=-\frac{V_0}{R}\sqrt{\frac{\pi}{3}}
\]


\subsection{Potencial en el interior}


Sustituyendo en (1) (solo sobreviven $\ell=0$ y $\ell=1$, $m=0$):

\[
\phi(r,\theta)=A_{00}\,Y_0^0+A_{10}\,r\,Y_1^0=V_0\sqrt{\pi}\cdot\frac{1}{\sqrt{4\pi}}-\frac{V_0}{R}\sqrt{\frac{\pi}{3}}\cdot r\cdot\sqrt{\frac{3}{4\pi}}\cos\theta
\]


\[
\boxed{\phi(r,\theta)=\frac{V_0}{2}\left(1-\frac{r}{R}\cos\theta\right)}\qquad (r\le R)\tag{3}
\]


\subsection{Campo eléctrico}


\[
\vec E=-\nabla\phi=-\left[\hat r\,\frac{\partial}{\partial r}+\frac{\hat\theta}{r}\,\frac{\partial}{\partial\theta}\right]\phi
\]


Con $\dfrac{\partial\phi}{\partial r}=-\dfrac{V_0}{2R}\cos\theta$ y $\dfrac{\partial\phi}{\partial\theta}=\dfrac{V_0\,r}{2R}\sin\theta$, se obtiene:

\[
\boxed{\vec E=\frac{V_0}{2R}\left(\cos\theta\,\hat r-\sin\theta\,\hat\theta\right)}\qquad (r<R)\tag{4}
\]


\begin{tcolorbox}[colback=blue!3, colframe=blue!35!black, title={Nota}, fonttitle=\bfseries]
\textbf{Nota:} $\cos\theta\,\hat r-\sin\theta\,\hat\theta=\hat z$ es una identidad válida en cualquier punto del espacio. Por lo tanto:

\[
\vec E=\frac{V_0}{2R}\,\hat z
\]

¡El campo en el \textbf{interior es uniforme}, aunque el potencial impuesto en la frontera, $V_0\sin^2(\theta/2)$, no lo parezca a simple vista!

\end{tcolorbox}


\medskip


\section{Teorema de adición de armónicos esféricos}


Sean $(\theta,\varphi)$ y $(\theta',\varphi')$ dos direcciones arbitrarias, y sea $\gamma$ el ángulo entre ellas (el ángulo entre los vectores unitarios $\hat r$ y $\hat r\,'$). El \textbf{teorema de adición} establece:

\[
P_\ell(\cos\gamma)=\frac{4\pi}{2\ell+1}\sum_{m=-\ell}^{\ell}(Y_\ell^m)^*(\theta,\varphi)\,Y_\ell^m(\theta',\varphi')
\]


\textbf{¿Por qué es importante?} Permite escribir $P_\ell(\cos\gamma)$ ---que depende del ángulo \textit{relativo} entre $\vec r$ y $\vec r\,'$--- como una suma de productos de armónicos esféricos evaluados \textit{por separado} en $(\theta,\varphi)$ y en $(\theta',\varphi')$. Esto es indispensable para integrar sobre una distribución de carga (variable $\vec r\,'$) manteniendo el punto de observación $\vec r$ arbitrario. Se utiliza en la sección 5.

\begin{center}
\begin{tikzpicture}[>=Stealth]
  \node[draw=black!65, fill=green!6, rounded corners=2pt, very thick,
        minimum width=0.82\textwidth, minimum height=2.0cm, align=center] (board) {};
  \node[align=center] at (board.center) {\Large\itshape ``Estudiar, continuar el ejercicio''\\[2mm]
  \normalsize Teorema de adición: $P_\ell(\cos\gamma)\longrightarrow Y_\ell^m(\theta,\varphi)Y_\ell^m(\theta',\varphi')$};
  \draw[decorate, decoration={coil, aspect=0.25, segment length=3pt, amplitude=2pt}, black!45]
    ([xshift=-3.6cm,yshift=-0.65cm]board.center) -- ([xshift=3.6cm,yshift=-0.65cm]board.center);
\end{tikzpicture}
\end{center}


\medskip


\section{Función de Green para la ecuación de Poisson}


\[
\nabla^2\phi=-\frac{\rho}{\varepsilon_0}\qquad\text{(Ecuación de Poisson)}
\]


La función de Green del laplaciano, $G(\vec r,\vec r\,')=\dfrac{1}{4\pi|\vec r-\vec r\,'|}$, satisface $\nabla^2G=-\delta^{3}(\vec r-\vec r\,')$, de modo que la solución general de Poisson es:

\[
\boxed{\phi(\vec r)=\frac{1}{4\pi\varepsilon_0}\int\frac{\rho(\vec r\,')}{|\vec r-\vec r\,'|}\,dV'}\tag{5}
\]


(esta es, simplemente, la generalización de la Ley de Coulomb a una distribución continua de carga $\rho(\vec r\,')$).

\subsection{\texorpdfstring{Expansión multipolar de $1/|\vec r-\vec r\,'|$}{Expansión multipolar del potencial de Green}}


Por la ley de cosenos, $|\vec r-\vec r\,'|=\sqrt{r^2+r'^2-2rr'\cos\gamma}$.

Caso $r>r'$:

\[
\frac{1}{|\vec r-\vec r\,'|}=\frac{1}{r}\cdot\frac{1}{\sqrt{1+(r'/r)^2-2(r'/r)\cos\gamma}}
\]


Con $t=r'/r$ y $x=\cos\gamma$, se reconoce la \textbf{función generatriz de los polinomios de Legendre}:

\[
g(t,x)=\frac{1}{\sqrt{1-2xt+t^2}}=\sum_{\ell=0}^{\infty}P_\ell(x)\,t^{\ell}\qquad(|t|<1)
\]


por lo que

\[
\frac{1}{|\vec r-\vec r\,'|}=\frac{1}{r}\sum_{\ell=0}^{\infty}P_\ell(\cos\gamma)\left(\frac{r'}{r}\right)^{\ell}\qquad(r>r')
\]


Caso $r<r'$ (intercambiando $r\leftrightarrow r'$ por simetría):

\[
\frac{1}{|\vec r-\vec r\,'|}=\frac{1}{r'}\sum_{\ell=0}^{\infty}P_\ell(\cos\gamma)\left(\frac{r}{r'}\right)^{\ell}\qquad(r<r')
\]


\textbf{Forma unificada}, definiendo $r_>=\max(r,r')$ y $r_<=\min(r,r')$:

\[
\boxed{\frac{1}{|\vec r-\vec r\,'|}=\sum_{\ell=0}^{\infty}\frac{r_<^{\,\ell}}{r_>^{\,\ell+1}}\,P_\ell(\cos\gamma)}\tag{6}
\]


Esta es la célebre \textbf{expansión multipolar}, base del siguiente ejercicio.

\begin{center}
\begin{tikzpicture}[scale=1.15, >=Stealth]
  \coordinate (O) at (0,0);
  \coordinate (R) at (28:3.4);
  \coordinate (Rp) at (72:2.4);
  \draw[->, thick] (O) -- (3.9,0) node[right] {$x$};
  \draw[->, thick] (O) -- (0,3.0) node[above] {$z$};
  \draw[->, very thick, blue!70!black] (O) -- (R) node[midway, below right] {$\vec r$};
  \draw[->, very thick, red!70!black] (O) -- (Rp) node[midway, left] {$\vec r\,'$};
  \draw[dashed, thick, purple!70!black] (Rp) -- (R) node[midway, above right] {$\vec r-\vec r\,'$};
  \pic[draw=black, ->, "$\gamma$", angle radius=0.75cm, angle eccentricity=1.25] {angle=R--O--Rp};
  \fill (O) circle (1.3pt) node[below left] {$O$};
  \fill (R) circle (1.4pt) node[right] {observación};
  \fill (Rp) circle (1.4pt) node[above] {fuente};
\end{tikzpicture}
\end{center}


\medskip


\section{\texorpdfstring{Ejercicio: esfera maciza con carga $+Q$ uniforme (aislada)}{Ejercicio: esfera maciza con carga positiva uniforme}}


\textbf{Enunciado:} calcular el potencial $\phi(\vec r)$ y el campo $\vec E(\vec r)$ generados por una esfera maciza de radio $R$, con densidad de carga volumétrica \textbf{uniforme} y carga total $+Q$, aislada (sin influencias externas).

\[
\rho_0=\frac{Q}{\frac{4}{3}\pi R^3}=\frac{3Q}{4\pi R^3}\quad(\text{constante})
\]


\begin{center}
\begin{tikzpicture}[scale=0.95, >=Stealth]
  \coordinate (C) at (0,0);
  \coordinate (S) at (48:1.45);
  \coordinate (P) at (22:3.25);
  \draw[fill=orange!10, draw=orange!70!black, thick] (C) circle (1.85);
  \draw[->, thick] (C) -- (0,2.4) node[above] {$z$};
  \draw[->, thick] (C) -- (2.35,0) node[right] {$x$};
  \draw[->, very thick, blue!70!black] (C) -- (P) node[midway, below right] {$\vec r$};
  \draw[->, very thick, red!70!black] (C) -- (S) node[midway, left] {$\vec r\,'$};
  \draw[dashed, thick, purple!70!black] (S) -- (P) node[midway, above] {$\vec r-\vec r\,'$};
  \draw[<->, orange!80!black] (C) -- (-1.85,0) node[midway, below] {$R$};
  \pic[draw=black, ->, "$\gamma$", angle radius=0.58cm, angle eccentricity=1.3] {angle=P--C--S};
  \fill (C) circle (1.4pt) node[below left] {$C$};
  \fill (S) circle (1.4pt) node[above left] {fuente};
  \fill (P) circle (1.4pt) node[right] {$P$};
  \node[align=center] at (0,-2.35) {$\rho_0=\mathrm{cte}$\\carga total $+Q$};
  \begin{scope}[xshift=6.0cm, yshift=-0.35cm]
    \coordinate (O) at (0,0);
    \coordinate (Zs) at (0,2.1);
    \draw[->, thick] (O) -- (0,2.1) node[above] {$z$};
    \draw[->, thick] (O) -- (2.2,0) node[right] {$x$};
    \draw[->, thick] (O) -- (-1.35,-1.05) node[below left] {$y$};
    \draw[->, very thick, blue!70!black] (O) -- (45:1.75) node[midway, above left] {$r$};
    \draw[dashed] (45:1.75) -- (1.24,0);
    \coordinate (Ps) at (45:1.75);
    \pic[draw=black, ->, "$\theta$", angle radius=0.55cm, angle eccentricity=1.35] {angle=Ps--O--Zs};
    \draw[->] (0.55,0) arc[start angle=0, end angle=-36, radius=0.55] node[midway, right] {$\varphi$};
    \node at (0,-1.65) {ejes esféricos};
  \end{scope}
\end{tikzpicture}
\end{center}


\subsection{Estrategia}


Se combina la fórmula de Green (5) con la expansión multipolar (6):

\[
\phi(\vec r)=\frac{\rho_0}{4\pi\varepsilon_0}\sum_{\ell=0}^{\infty}\int_0^{R}\!\int_{4\pi}\frac{r_<^{\,\ell}}{r_>^{\,\ell+1}}\,P_\ell(\cos\gamma)\;r'^2\,dr'\,d\Omega'
\]


\textbf{Simplificación clave:} como $\rho_0$ es constante (no depende de $\theta',\varphi'$), al integrar $P_\ell(\cos\gamma)$ sobre $d\Omega'$ ---usando el teorema de adición de la sección 3--- \textbf{todos los términos con $\ell>0$ se anulan}. Físicamente: una esfera uniformemente cargada no tiene momentos multipolares de orden $\ell\geq1$ respecto a su centro. Solo sobrevive $\ell=0$, con $P_0=1$ y $\int d\Omega'=4\pi$.

\subsection{\texorpdfstring{Región exterior ($r>R\;\Rightarrow\; r_>=r,\ r_<=r'$ siempre)}{Región exterior}}


\[
\phi(r)=\frac{\rho_0}{4\pi\varepsilon_0}\cdot\frac{4\pi}{r}\int_0^{R} r'^2\,dr'=\frac{\rho_0}{\varepsilon_0\,r}\cdot\frac{R^3}{3}=\frac{Q}{4\pi\varepsilon_0\,r}=\frac{kQ}{r}
\]


\subsection{\texorpdfstring{Región interior ($r\le R$)}{Región interior}}


Aquí hay que separar la integral en $r'<r$ (donde $r_>=r$) y $r'>r$ (donde $r_>=r'$):

\[
\phi(r)=\frac{\rho_0}{\varepsilon_0}\left[\frac{1}{r}\int_0^{r}r'^2\,dr'+\int_{r}^{R}r'\,dr'\right]=\frac{\rho_0}{2\varepsilon_0}\left(R^2-\frac{r^2}{3}\right)=\frac{kQ}{2R^3}\left(3R^2-r^2\right)
\]


\subsection{Resultados finales}


\[
\boxed{\phi(\vec r)=\begin{cases}\dfrac{kQ}{r}\,, & r>R\\[3mm] \dfrac{kQ}{2R^3}\left(3R^2-r^2\right)\,, & r\le R\end{cases}}\qquad\left(k\equiv\dfrac{1}{4\pi\varepsilon_0}\right)
\]


Por simetría esférica, $\vec E=-\dfrac{\partial\phi}{\partial r}\,\hat r$, de donde:

\[
\boxed{\vec E(\vec r)=\begin{cases}\dfrac{1}{4\pi\varepsilon_0}\dfrac{Q}{r^2}\,\hat r\,, & r>R\\[3mm] \dfrac{kQ\,r}{R^3}\,\hat r\,, & r\le R\end{cases}}
\]


Este es el resultado clásico para una esfera uniformemente cargada: fuera se comporta como una carga puntual $Q$ en el origen, y dentro el campo crece linealmente con $r$.

\medskip


\section{Ecuación de Laguerre}


\begin{tcolorbox}[colback=blue!3, colframe=blue!35!black, title={Nota}, fonttitle=\bfseries]
\textbf{Contexto:} esta ecuación aparece al resolver la parte \textbf{radial} de la ecuación de Schrödinger para el átomo de hidrógeno (potencial coulombiano $\sim 1/r$). Así como los armónicos esféricos $Y_\ell^m(\theta,\varphi)$ resuelven la parte angular (secciones 1–5), los \textbf{polinomios de Laguerre asociados} resuelven la parte radial.

\end{tcolorbox}


\subsection{Ecuación de Laguerre (forma simple)}


\[
x\ddot y+(1-x)\dot y+ny=0\tag{7}
\]


(donde $\dot y=dy/dx$). Se resuelve por el \textbf{método de Frobenius} (serie de potencias alrededor del punto singular regular $x=0$). La solución polinómica, para $n=0,1,2,\dots$, es el \textbf{polinomio de Laguerre}:

\[
y=L_n(x)=\sum_{p=0}^{n}\frac{(-1)^p\,x^p}{(p!)^2\,(n-p)!}\tag{8}
\]


\textbf{Fórmula de Rodrigues} (forma compacta equivalente):

\[
L_n(x)=\frac{e^{x}}{n!}\,\frac{d^{\,n}}{dx^{n}}\left(x^{n}e^{-x}\right)\tag{9}
\]


\textbf{Relación de ortogonalidad} en $[0,\infty)$ con peso $e^{-x}$:

\[
\int_0^\infty e^{-x}\,L_n(x)\,L_m(x)\,dx=\delta_{nm}
\]


\subsection{Ecuación general (asociada) de Laguerre}


\[
x\ddot y+(k+1-x)\dot y+ny=0\tag{10}
\]


\textbf{Solución} --- polinomios de Laguerre asociados (o generalizados):

\[
\boxed{L_n^{k}(x)=\frac{e^{x}\,x^{-k}}{n!}\,\frac{d^{\,n}}{dx^{n}}\left(x^{\,n+k}\,e^{-x}\right)}\tag{11}
\]


Caso particular: para $k=0$ se recupera $L_n^{0}(x)=L_n(x)$ (los polinomios "simples" de la sección 6.1).

\begin{tcolorbox}[colback=blue!3, colframe=blue!35!black, title={Nota}, fonttitle=\bfseries]
\textbf{Nota --- analogía con Legendre:} así como los polinomios de Legendre $P_\ell(x)$ se generalizan a los \textbf{asociados} $P_\ell^{m}(x)$ (presentes en $Y_\ell^m$), los polinomios de Laguerre $L_n(x)$ se generalizan a los \textbf{asociados} $L_n^{k}(x)$. En el átomo de hidrógeno, la función de onda completa tiene la forma

\[
\psi_{n\ell m}(r,\theta,\varphi)=R_{n\ell}(r)\,Y_\ell^m(\theta,\varphi),
\]

donde la parte radial $R_{n\ell}(r)$ se construye con $L_{n}^{k}(x)$ (con $k=2\ell+1$ y $x\propto r$), mientras que la parte angular $Y_\ell^m$ involucra $P_\ell^{m}(\cos\theta)$ --- cerrando así el círculo con la sección 1.

\end{tcolorbox}


\begin{center}
\begin{tikzpicture}[scale=1.0, >=Stealth]
  \coordinate (O) at (0,0);
  \coordinate (P) at (42:2.5);
  \coordinate (Z) at (0,3.0);
  \draw[->, thick] (O) -- (0,3.0) node[above] {$z$};
  \draw[->, thick] (O) -- (3.2,0) node[right] {$x$};
  \draw[->, thick] (O) -- (-1.8,-1.25) node[below left] {$y$};
  \draw[->, very thick, blue!70!black] (O) -- (P) node[midway, above left] {$r$};
  \draw[dashed] (P) -- (1.86,0);
  \draw[dashed] (P) -- (-0.55,-0.38);
  \pic[draw=black, ->, "$\theta$", angle radius=0.75cm, angle eccentricity=1.25] {angle=P--O--Z};
  \draw[->] (0.85,0) arc[start angle=0, end angle=-34, radius=0.85] node[midway, right] {$\varphi$};
  \fill (O) circle (1.3pt) node[below left] {$O$};
  \fill (P) circle (1.3pt) node[above right] {$\psi_{n\ell m}=R_{n\ell}(r)Y_\ell^m(\theta,\varphi)$};
  \node[draw=black!45, fill=blue!4, rounded corners=2pt, align=center] at (5.3,1.0)
    {parte radial\\$R_{n\ell}(r)$\\Laguerre};
  \node[draw=black!45, fill=green!5, rounded corners=2pt, align=center] at (5.3,-1.0)
    {parte angular\\$Y_\ell^m(\theta,\varphi)$\\Legendre asociado};
  \draw[->, thick] (3.6,0.45) -- (4.35,0.85);
  \draw[->, thick] (3.6,-0.2) -- (4.35,-0.85);
\end{tikzpicture}
\end{center}

\section{Problema 1: Determinar $A_{\ell m}$}
Demostrar que los coeficientes $A_{\ell m}$ en la expansión de una función arbitraria $f(\theta, \phi)$ en términos de armónicos esféricos:
$$f(\theta, \phi) = \sum_{\ell=0}^{\infty} \sum_{m=-\ell}^{\ell} A_{\ell m} Y_\ell^m(\theta, \phi)$$
están dados por la expresión:
$$A_{\ell m} = \int_{4\pi} f(\theta, \phi) (Y_\ell^m(\theta, \phi))^* \, d\Omega$$

\subsection{Solución}

Partimos de la expansión de la función $f(\theta, \phi)$ en la base completa de armónicos esféricos. Para evitar confusión con los índices específicos $\ell$ y $m$ que deseamos aislar en el lado izquierdo de la ecuación final, reescribimos la sumatoria utilizando índices mudos $\ell'$ y $m'$:
$$f(\theta, \phi) = \sum_{\ell'=0}^{\infty} \sum_{m'=-\ell'}^{\ell'} A_{\ell' m'} Y_{\ell'}^{m'}(\theta, \phi)$$

Multiplicamos ambos lados de la ecuación anterior por el complejo conjugado del armónico esférico $(Y_\ell^m(\theta, \phi))^*$, correspondiente a los índices $\ell$ y $m$ que queremos proyectar:
$$f(\theta, \phi) (Y_\ell^m(\theta, \phi))^* = \left[ \sum_{\ell'=0}^{\infty} \sum_{m'=-\ell'}^{\ell'} A_{\ell' m'} Y_{\ell'}^{m'}(\theta, \phi) \right] (Y_\ell^m(\theta, \phi))^*$$

Integramos ambos lados de la ecuación sobre todo el ángulo sólido $4\pi$ (es decir, $\theta$ desde $0$ hasta $\pi$ y $\phi$ desde $0$ hasta $2\pi$, con el elemento diferencial de ángulo sólido $d\Omega = \sen\theta \, d\theta \, d\phi$):
$$\int_{4\pi} f(\theta, \phi) (Y_\ell^m(\theta, \phi))^* \, d\Omega = \int_{4\pi} \left[ \sum_{\ell'=0}^{\infty} \sum_{m'=-\ell'}^{\ell'} A_{\ell' m'} Y_{\ell'}^{m'}(\theta, \phi) (Y_\ell^m(\theta, \phi))^* \right] d\Omega$$

Debido a la linealidad del operador integral, podemos intercambiar el orden de la integración y las sumatorias discretas, sacando los coeficientes constantes $A_{\ell' m'}$ fuera de la integral:
$$\int_{4\pi} f(\theta, \phi) (Y_\ell^m(\theta, \phi))^* \, d\Omega = \sum_{\ell'=0}^{\infty} \sum_{m'=-\ell'}^{\ell'} A_{\ell' m'} \left[ \int_{4\pi} Y_{\ell'}^{m'}(\theta, \phi) (Y_\ell^m(\theta, \phi))^* \, d\Omega \right]$$

Aquí utilizamos la propiedad fundamental de ortonormalidad de los armónicos esféricos, dada por la ecuación:
$$\int_{4\pi} (Y_\ell^m(\theta, \phi))^* Y_{\ell'}^{m'}(\theta, \phi) \, d\Omega = \delta_{\ell \ell'} \delta_{m m'}$$

Sustituimos la condición de ortonormalidad en la expresión del Paso 4:
$$\int_{4\pi} f(\theta, \phi) (Y_\ell^m(\theta, \phi))^* \, d\Omega = \sum_{\ell'=0}^{\infty} \sum_{m'=-\ell'}^{\ell'} A_{\ell' m'} \delta_{\ell \ell'} \delta_{m m'}$$
La presencia de las deltas de Kronecker $\delta_{\ell \ell'} \delta_{m m'}$ actúa como un filtro que anula todos los términos de la doble sumatoria, excepto aquel en el que $\ell' = \ell$ y $m' = m$. Por lo tanto, la sumatoria colapsa a un único término:
$$\int_{4\pi} f(\theta, \phi) (Y_\ell^m(\theta, \phi))^* \, d\Omega = A_{\ell m}$$

Reordenando la igualdad, obtenemos exactamente la expresión que queríamos demostrar:
$$A_{\ell m} = \int_{4\pi} f(\theta, \phi) (Y_\ell^m(\theta, \phi))^* \, d\Omega$$



\section{Problema 2: $\vec{r} \cdot \vec{\nabla} Y_\ell^m(\theta,\phi) = 0$}

Demostrar que el producto punto entre el vector de posición y el gradiente de un armónico esférico es nulo:
$$\vec{r} \cdot \vec{\nabla} Y_\ell^m(\theta,\phi) = 0$$

\subsection{Solución}

En coordenadas esféricas $(r, \theta, \phi)$, el operador gradiente se define como:
$$\vec{\nabla} = \hat{e}_r \frac{\partial}{\partial r} + \hat{e}_\theta \frac{1}{r} \frac{\partial}{\partial \theta} + \hat{e}_\phi \frac{1}{r \sen\theta} \frac{\partial}{\partial \phi}$$
El vector de posición $\vec{r}$ en este sistema de coordenadas tiene una expresión extremadamente simple, apuntando puramente en la dirección radial con magnitud $r$:
$$\vec{r} = r \hat{e}_r$$
Al realizar el producto punto entre $\vec{r}$ y el operador $\vec{\nabla}$, aprovechamos la ortonormalidad de la base local ($\hat{e}_r \cdot \hat{e}_r = 1$, $\hat{e}_r \cdot \hat{e}_\theta = 0$, $\hat{e}_r \cdot \hat{e}_\phi = 0$):
$$\vec{r} \cdot \vec{\nabla} = (r \hat{e}_r) \cdot \left( \hat{e}_r \frac{\partial}{\partial r} + \hat{e}_\theta \frac{1}{r} \frac{\partial}{\partial \theta} + \hat{e}_\phi \frac{1}{r \sen\theta} \frac{\partial}{\partial \phi} \right) = r \frac{\partial}{\partial r}$$
Este operador $r \frac{\partial}{\partial r}$ mide la tasa de cambio de una función a lo largo de la dirección radial (es decir, cómo cambia la función si nos alejamos o acercamos al origen, manteniendo $\theta$ y $\phi$ constantes).

Por definición, los armónicos esféricos $Y_\ell^m(\theta,\phi)$ son funciones que dependen \textbf{exclusivamente} de las coordenadas angulares $\theta$ y $\phi$. No poseen dependencia alguna con la coordenada radial $r$. Matemáticamente, esto se expresa como:
$$\frac{\partial}{\partial r} Y_\ell^m(\theta,\phi) = 0$$
Por lo tanto, al aplicar el operador $\vec{r} \cdot \vec{\nabla}$ sobre el armónico esférico, obtenemos:
$$\vec{r} \cdot \vec{\nabla} Y_\ell^m(\theta,\phi) = r \frac{\partial}{\partial r} Y_\ell^m(\theta,\phi) = r (0) = 0$$

\section{Problema 3: Cascarones esféricos concentricos}

Si en la región entre dos cascarones esféricos de radios $a$ y $b$ se satisface la ecuación de Laplace para potencial electrostático y si $\phi = V_1(\theta, \varphi)$ en $r = a$ y $\phi = V_2(\theta, \varphi)$ en $r = b$, evaluar $\phi(r, \theta, \varphi)$.

Obsérvese que $r = 0$ no está incluido en el problema, de modo que $B_{\ell m}$ será en principio diferente de cero.

\subsection{Solución}

Dado que el potencial satisface la ecuación de Laplace en la región $a < r < b$, y considerando que el origen no está incluido en el dominio del problema, la solución general en coordenadas esféricas se expresa como una combinación lineal de armónicos esféricos:

$$
\phi(r,\theta,\varphi) = \sum_{\ell=0}^{\infty} \sum_{m=-\ell}^{\ell} \left( A_{\ell m} r^\ell + \frac{B_{\ell m}}{r^{\ell+1}} \right) Y_\ell^m(\theta,\varphi)
$$

donde $Y_\ell^m(\theta,\varphi)$ son los armónicos esféricos, y tanto los coeficientes $A_{\ell m}$ como $B_{\ell m}$ deben determinarse a partir de las condiciones de frontera.

Las condiciones de frontera son:
$$
\phi(a,\theta,\varphi) = V_1(\theta,\varphi)
$$
$$
\phi(b,\theta,\varphi) = V_2(\theta,\varphi)
$$

Sustituyendo $r=a$ y $r=b$ en la solución general:
$$
\sum_{\ell,m} \left( A_{\ell m} a^\ell + \frac{B_{\ell m}}{a^{\ell+1}} \right) Y_\ell^m(\theta,\varphi) = V_1(\theta,\varphi)
$$
$$
\sum_{\ell,m} \left( A_{\ell m} b^\ell + \frac{B_{\ell m}}{b^{\ell+1}} \right) Y_\ell^m(\theta,\varphi) = V_2(\theta,\varphi)
$$

Utilizando la ortogonalidad de los armónicos esféricos:
$$
\int \left(Y_{\ell'}^{m'}(\theta,\varphi)\right)^* Y_\ell^m(\theta,\varphi) d\Omega = \delta_{\ell\ell'}\delta_{mm'}
$$

Multiplicamos ambas ecuaciones por $\left(Y_{\ell'}^{m'}(\theta,\varphi)\right)^*$ e integramos sobre el ángulo sólido:
$$
A_{\ell m} a^\ell + \frac{B_{\ell m}}{a^{\ell+1}} = \int V_1(\theta,\varphi) \left(Y_\ell^m(\theta,\varphi)\right)^* d\Omega \equiv C_{\ell m}
$$
$$
A_{\ell m} b^\ell + \frac{B_{\ell m}}{b^{\ell+1}} = \int V_2(\theta,\varphi) \left(Y_\ell^m(\theta,\varphi)\right)^* d\Omega \equiv D_{\ell m}
$$

Esto nos da un sistema de dos ecuaciones con dos incógnitas ($A_{\ell m}$ y $B_{\ell m}$) para cada par $(\ell,m)$:
$$
A_{\ell m} a^\ell + B_{\ell m} a^{-(\ell+1)} = C_{\ell m}
$$
$$
A_{\ell m} b^\ell + B_{\ell m} b^{-(\ell+1)} = D_{\ell m}
$$

Resolviendo este sistema:
$$
A_{\ell m} = \frac{b^{\ell+1}C_{\ell m} - a^{\ell+1}D_{\ell m}}{b^{2\ell+1} - a^{2\ell+1}}, \qquad 
B_{\ell m} = \frac{a^\ell b^\ell (D_{\ell m} - C_{\ell m})}{b^{-(\ell+1)} - a^{-(\ell+1)}} = \frac{a^\ell b^\ell (C_{\ell m} - D_{\ell m})}{a^{-(\ell+1)} - b^{-(\ell+1)}}
$$

Por lo tanto, el potencial electrostático en la región entre los dos cascarones es:
$$
\phi(r,\theta,\varphi) = \sum_{\ell=0}^{\infty} \sum_{m=-\ell}^{\ell} \left[ \frac{b^{\ell+1}C_{\ell m} - a^{\ell+1}D_{\ell m}}{b^{2\ell+1} - a^{2\ell+1}} r^\ell + \frac{a^\ell b^\ell (C_{\ell m} - D_{\ell m})}{a^{-(\ell+1)} - b^{-(\ell+1)}} \frac{1}{r^{\ell+1}} \right] Y_\ell^m(\theta,\varphi)
$$

donde los coeficientes $C_{\ell m}$ y $D_{\ell m}$ están dados por las proyecciones de los potenciales de frontera sobre los armónicos esféricos.

\section{Problema 4: Temperatura sobre superficie esférica}
La temperatura sobre la superficie de una esfera de radio $R$ se mantiene fija y con un valor $T(R, \theta, \varphi) = f(\theta)$. Demuestre que en esta situación de estado estacionario:
$$
T(r,\theta,\varphi) = \frac{1}{2} \sum_{\ell=0}^{\infty} \frac{(2\ell+1)}{R^\ell} \left[ \int_0^{\pi} f(\theta') \sen\theta' P_\ell(\cos\theta') d\theta' \right] r^\ell P_\ell(\cos\theta)
$$

\subsection{Solución}

En estado estacionario, la temperatura satisface la ecuación de Laplace:
$$
\nabla^2 T(r,\theta,\varphi) = 0
$$

Dado que la condición de frontera $T(R,\theta,\varphi) = f(\theta)$ es independiente del ángulo azimutal $\varphi$, la solución también será independiente de $\varphi$ (simetría azimutal). Por lo tanto, $T = T(r,\theta)$.

La solución general de la ecuación de Laplace con simetría azimutal en coordenadas esféricas se expresa en términos de los polinomios de Legendre:
$$
T(r,\theta) = \sum_{\ell=0}^{\infty} \left( A_\ell r^\ell + \frac{B_\ell}{r^{\ell+1}} \right) P_\ell(\cos\theta)
$$

Como estamos interesados en la región interior de la esfera ($r < R$) y la temperatura debe ser finita en el origen ($r=0$), debemos tener $B_\ell = 0$ para todo $\ell$. Así:
$$
T(r,\theta) = \sum_{\ell=0}^{\infty} A_\ell r^\ell P_\ell(\cos\theta)
$$

Aplicando la condición de frontera en $r = R$:
$$
T(R,\theta) = f(\theta) = \sum_{\ell=0}^{\infty} A_\ell R^\ell P_\ell(\cos\theta)
$$

Para determinar los coeficientes $A_\ell$, utilizamos la ortogonalidad de los polinomios de Legendre:
$$
\int_0^{\pi} P_{\ell'}(\cos\theta) P_\ell(\cos\theta) \sen\theta \, d\theta = \frac{2}{2\ell+1} \delta_{\ell\ell'}
$$

Multiplicamos la ecuación de la condición de frontera por $P_{\ell'}(\cos\theta) \sen\theta$ e integramos de $0$ a $\pi$:
$$
\int_0^{\pi} f(\theta) P_\ell(\cos\theta) \sen\theta \, d\theta = \sum_{\ell'=0}^{\infty} A_{\ell'} R^{\ell'} \int_0^{\pi} P_{\ell'}(\cos\theta) P_\ell(\cos\theta) \sen\theta \, d\theta
$$

Utilizando la relación de ortogonalidad:
$$
\int_0^{\pi} f(\theta) P_\ell(\cos\theta) \sen\theta \, d\theta = A_\ell R^\ell \cdot \frac{2}{2\ell+1}
$$

Despejando $A_\ell$:
$$
A_\ell = \frac{2\ell+1}{2R^\ell} \int_0^{\pi} f(\theta) P_\ell(\cos\theta) \sen\theta \, d\theta
$$

Sustituyendo este resultado en la expresión para $T(r,\theta)$:
$$
T(r,\theta) = \sum_{\ell=0}^{\infty} \left[ \frac{2\ell+1}{2R^\ell} \int_0^{\pi} f(\theta') P_\ell(\cos\theta') \sen\theta' \, d\theta' \right] r^\ell P_\ell(\cos\theta)
$$

Reorganizando los términos, obtenemos la expresión solicitada:
$$
T(r,\theta,\varphi) = \frac{1}{2} \sum_{\ell=0}^{\infty} \frac{(2\ell+1)}{R^\ell} \left[ \int_0^{\pi} f(\theta') \sen\theta' P_\ell(\cos\theta') d\theta' \right] r^\ell P_\ell(\cos\theta)
$$

que es independiente de $\varphi$ debido a la simetría azimutal del problema.

\section{Problema 5: Ecuación de Helmholtz}
Demuestre que la solución a la ecuación de Helmholtz $(\nabla^2 + k^2)\phi = 0$, en coordenadas esféricas, que satisface la condición de continuidad en $\varphi$, e incluye $\theta = 0, \pi/2$, es:
$$\phi(r, \theta, \varphi) = \sum_{\ell=0}^{\infty} \sum_{m=-\ell}^{\ell} [A_{\ell m} j_\ell(kr) + B_{\ell m} \eta_\ell(kr)] Y_\ell^m(\theta, \varphi)$$

\subsection{Solución}
La ecuación de Helmholtz es:
$$(\nabla^2 + k^2)\phi(r, \theta, \varphi) = 0$$

En coordenadas esféricas, el laplaciano tiene la forma:
$$\nabla^2\phi = \frac{1}{r^2}\frac{\partial}{\partial r}\left(r^2\frac{\partial\phi}{\partial r}\right) + \frac{1}{r^2\sin\theta}\frac{\partial}{\partial\theta}\left(\sin\theta\frac{\partial\phi}{\partial\theta}\right) + \frac{1}{r^2\sin^2\theta}\frac{\partial^2\phi}{\partial\varphi^2}$$

Proponemos una solución separada de la forma:
$$\phi(r, \theta, \varphi) = R(r)\Theta(\theta)\Phi(\varphi)$$

Sustituyendo en la ecuación de Helmholtz y dividiendo por $\phi = R\Theta\Phi$:
$$\frac{1}{R}\frac{1}{r^2}\frac{d}{dr}\left(r^2\frac{dR}{dr}\right) + \frac{1}{\Theta}\frac{1}{r^2\sin\theta}\frac{d}{d\theta}\left(\sin\theta\frac{d\Theta}{d\theta}\right) + \frac{1}{\Phi}\frac{1}{r^2\sin^2\theta}\frac{d^2\Phi}{d\varphi^2} + k^2 = 0$$

Multiplicando por $r^2$:
$$\frac{1}{R}\frac{d}{dr}\left(r^2\frac{dR}{dr}\right) + \frac{1}{\Theta}\frac{1}{\sin\theta}\frac{d}{d\theta}\left(\sin\theta\frac{d\Theta}{d\theta}\right) + \frac{1}{\Phi}\frac{1}{\sin^2\theta}\frac{d^2\Phi}{d\varphi^2} + k^2r^2 = 0$$

El término que depende de $\varphi$ debe ser constante. Sea $m^2$ esta constante:
$$\frac{1}{\Phi}\frac{d^2\Phi}{d\varphi^2} = -m^2$$

La solución es $\Phi(\varphi) = e^{im\varphi}$. La condición de continuidad $\Phi(\varphi) = \Phi(\varphi + 2\pi)$ exige que $m$ sea entero: $m = 0, \pm 1, \pm 2, \ldots$

Separando la parte angular, obtenemos:
$$\frac{1}{\Theta}\frac{1}{\sin\theta}\frac{d}{d\theta}\left(\sin\theta\frac{d\Theta}{d\theta}\right) - \frac{m^2}{\sin^2\theta} = -\ell(\ell+1)$$

donde $\ell(\ell+1)$ es la constante de separación. Esta es la \textbf{ecuación asociada de Legendre}:
$$\frac{1}{\sin\theta}\frac{d}{d\theta}\left(\sin\theta\frac{d\Theta}{d\theta}\right) + \left[\ell(\ell+1) - \frac{m^2}{\sin^2\theta}\right]\Theta = 0$$

Con el cambio $x = \cos\theta$, esta ecuación se transforma en:
$$\frac{d}{dx}\left[(1-x^2)\frac{dP}{dx}\right] + \left[\ell(\ell+1) - \frac{m^2}{1-x^2}\right]P = 0$$

Las soluciones regulares en $x = \pm 1$ (es decir, $\theta = 0, \pi$) son los \textbf{polinomios asociados de Legendre} $P_\ell^m(\cos\theta)$, donde $\ell$ debe ser entero no negativo y $|m| \leq \ell$.

Combinando las soluciones angulares, definimos los \textbf{armónicos esféricos}:
$$Y_\ell^m(\theta, \varphi) = N_{\ell m} P_\ell^m(\cos\theta) e^{im\varphi}$$

donde $N_{\ell m}$ es el factor de normalización. Los armónicos esféricos forman una base ortonormal completa en la esfera.

La parte radial satisface:
$$\frac{1}{r^2}\frac{d}{dr}\left(r^2\frac{dR}{dr}\right) + \left[k^2 - \frac{\ell(\ell+1)}{r^2}\right]R = 0$$

Haciendo el cambio $u(r) = rR(r)$, obtenemos:
$$\frac{d^2u}{dr^2} + \left[k^2 - \frac{\ell(\ell+1)}{r^2}\right]u = 0$$

Esta es la \textbf{ecuación de Bessel esférica}. Sus soluciones linealmente independientes son:
\begin{itemize}
\item $j_\ell(kr)$: función de Bessel esférica de primera especie (regular en $r=0$)
\item $\eta_\ell(kr)$: función de Bessel esférica de segunda especie (Neumann, singular en $r=0$)
\end{itemize}

Por lo tanto, la solución general radial es:
$$R_\ell(r) = A_\ell j_\ell(kr) + B_\ell \eta_\ell(kr)$$


Combinando todas las partes y sumando sobre todos los valores permitidos de $\ell$ y $m$, obtenemos:
$$\phi(r, \theta, \varphi) = \sum_{\ell=0}^{\infty} \sum_{m=-\ell}^{\ell} [A_{\ell m} j_\ell(kr) + B_{\ell m} \eta_\ell(kr)] Y_\ell^m(\theta, \varphi)$$

\section{Problema 6: Ecuación de onda}
Utilizando separación de variables para la ecuación de ondas en coordenadas esféricas demuestre que la solución tiene la forma:
$$\psi(\vec{r}, t) = \sum_{\ell=0}^{\infty} \sum_{m=-\ell}^{\ell} [A_{\ell m} j_\ell(kr) + B_{\ell m} \eta_\ell(kr)] Y_\ell^m(\theta, \phi) \times [C_{\ell m} e^{ikct} + D_{\ell m} e^{-ikct}]$$

\subsection{Solución}

La ecuación de ondas en tres dimensiones es:
$$\nabla^2\psi(\vec{r}, t) - \frac{1}{c^2}\frac{\partial^2\psi(\vec{r}, t)}{\partial t^2} = 0$$


Proponemos una solución separada:
$$\psi(\vec{r}, t) = \phi(\vec{r})T(t)$$

Sustituyendo en la ecuación de ondas:
$$T(t)\nabla^2\phi(\vec{r}) - \frac{1}{c^2}\phi(\vec{r})\frac{d^2T}{dt^2} = 0$$

Dividiendo por $\phi(\vec{r})T(t)$:
$$\frac{\nabla^2\phi}{\phi} = \frac{1}{c^2}\frac{1}{T}\frac{d^2T}{dt^2} = -k^2$$

donde $-k^2$ es la constante de separación (negativa para obtener soluciones oscilatorias).


La ecuación para $T(t)$ es:
$$\frac{d^2T}{dt^2} + k^2c^2T = 0$$

Cuya solución general es:
$$T(t) = C e^{ikct} + D e^{-ikct}$$

La parte espacial satisface la ecuación de Helmholtz:
$$(\nabla^2 + k^2)\phi(\vec{r}) = 0$$

Sabemos que la solución general en coordenadas esféricas es:
$$\phi(r, \theta, \varphi) = \sum_{\ell=0}^{\infty} \sum_{m=-\ell}^{\ell} [A_{\ell m} j_\ell(kr) + B_{\ell m} \eta_\ell(kr)] Y_\ell^m(\theta, \varphi)$$

Combinando las partes espacial y temporal, y permitiendo que los coeficientes temporales dependan de $\ell$ y $m$, obtenemos:
$$\psi(\vec{r}, t) = \sum_{\ell=0}^{\infty} \sum_{m=-\ell}^{\ell} [A_{\ell m} j_\ell(kr) + B_{\ell m} \eta_\ell(kr)] Y_\ell^m(\theta, \phi) \times [C_{\ell m} e^{ikct} + D_{\ell m} e^{-ikct}]$$

Esta es la solución general a la ecuación de ondas en coordenadas esféricas, expresada como superposición de modos normales caracterizados por los números cuánticos $\ell$ y $m$.

\section{Problema 7: Modos normales de oscilación}

Obtener el espectro de frecuencias (modos normales de oscilación) para una cavidad acústica resonante de forma esférica y radio $R$. La condición de frontera es de la forma: 

$$\frac{\partial\psi(\vec{r}, t)}{\partial r}\Bigg|_{r=R} = 0$$.

\subsection{Solución}

Para una cavidad acústica esférica, buscamos soluciones a la ecuación de ondas que satisfagan:
\begin{itemize}
\item La ecuación de ondas: $\nabla^2\psi - \frac{1}{c^2}\frac{\partial^2\psi}{\partial t^2} = 0$
\item Condición de frontera de Neumann: $\left.\frac{\partial\psi}{\partial r}\right|_{r=R} = 0$
\item Regularidad en el origen: $\psi$ debe ser finita en $r=0$
\end{itemize}


Sabemos que la solución general tiene la forma:
$$\psi(\vec{r}, t) = \sum_{\ell,m} [A_{\ell m} j_\ell(kr) + B_{\ell m} \eta_\ell(kr)] Y_\ell^m(\theta, \varphi) [C_{\ell m} e^{ikct} + D_{\ell m} e^{-ikct}]$$

La condición de regularidad en $r=0$ exige $B_{\ell m} = 0$ (pues $\eta_\ell(kr)$ diverge en el origen). Por lo tanto:
$$\psi(\vec{r}, t) = \sum_{\ell,m} A_{\ell m} j_\ell(kr) Y_\ell^m(\theta, \varphi) [C_{\ell m} e^{ikct} + D_{\ell m} e^{-ikct}]$$


La condición $\left.\frac{\partial\psi}{\partial r}\right|_{r=R} = 0$ debe cumplirse para cada modo $(\ell,m)$ individualmente. Esto implica:
$$\left.\frac{d}{dr}j_\ell(kr)\right|_{r=R} = 0$$

Usando la regla de la cadena:
$$k \cdot j_\ell'(kR) = 0$$

Como $k \neq 0$ (para soluciones no triviales), debemos tener:
$$j_\ell'(kR) = 0$$


Sea $\alpha_{\ell n}$ el $n$-ésimo cero de la derivada de la función de Bessel esférica $j_\ell'(x)$. Entonces:
$$k_{\ell n} R = \alpha_{\ell n}$$

de donde:
$$k_{\ell n} = \frac{\alpha_{\ell n}}{R}$$

La frecuencia angular correspondiente es:
$$\omega_{\ell n} = c k_{\ell n} = \frac{c \cdot \alpha_{\ell n}}{R}$$

y la frecuencia lineal:
$$f_{\ell n} = \frac{\omega_{\ell n}}{2\pi} = \frac{c \cdot \alpha_{\ell n}}{2\pi R}$$

Los modos normales de oscilación son:
$$\psi_{\ell n m}(\vec{r}, t) = j_\ell\left(\frac{\alpha_{\ell n} r}{R}\right) Y_\ell^m(\theta, \varphi) [C_{\ell n m} e^{i\omega_{\ell n}t} + D_{\ell n m} e^{-i\omega_{\ell n}t}]$$

con frecuencias:
$$\omega_{\ell n} = \frac{c \cdot \alpha_{\ell n}}{R}, \quad \ell = 0, 1, 2, \ldots; \quad n = 1, 2, 3, \ldots; \quad m = -\ell, \ldots, \ell$$

Cada frecuencia $\omega_{\ell n}$ es $(2\ell+1)$-veces degenerada debido a los diferentes valores de $m$.


\section{Problema 8: Partícula atómica en un cascarón esférico}

Una partícula atómica está confinada dentro de un cascarón esférico de radio $R$. La partícula es descrita por una función de onda que satisface la ecuación de Schrödinger


$$-\frac{\hbar^2}{2m}\nabla^2\psi = E\psi,$$


con la condición de que $\psi$ sea nula sobre las paredes. Demostrar que los niveles de energía permitidos tienen la forma:


$$E_{\ell n} = \frac{\alpha_{\ell n}^2\hbar^2}{2mR^2},$$


donde $\alpha_{\ell n}$ son los ceros de la función de Bessel esférica $j_\ell(x)$, es decir los ceros de $J_{\ell+1/2}(x)$, y que la función de onda es:


$$\psi = \sum_{\ell=0}^{\infty}\sum_{m=-\ell}^{\ell} A_{\ell m} j_\ell(\alpha_{\ell n}r/R) Y_{\ell}^m(\theta, \varphi).$$


\subsection{Solución}

La ecuación de Schrödinger independiente del tiempo para una partícula libre ($V=0$) en el interior del cascarón es:
$$-\frac{\hbar^2}{2m}\nabla^2\psi(r,\theta,\varphi) = E\psi(r,\theta,\varphi).$$

Reescribiendo:
$$\nabla^2\psi + k^2\psi = 0, \quad \text{donde} \quad k^2 = \frac{2mE}{\hbar^2}.$$

Proponemos una solución separada de la forma:
$$\psi(r,\theta,\varphi) = R(r)Y(\theta,\varphi).$$

Sustituyendo en la ecuación de Helmholtz y multiplicando por $r^2/\psi$:
$$\frac{1}{R}\frac{d}{dr}\left(r^2\frac{dR}{dr}\right) + k^2r^2 = -\frac{1}{Y}\left[\frac{1}{\sin\theta}\frac{\partial}{\partial\theta}\left(\sin\theta\frac{\partial Y}{\partial\theta}\right) + \frac{1}{\sin^2\theta}\frac{\partial^2 Y}{\partial\varphi^2}\right].$$

El lado izquierdo depende sólo de $r$ y el derecho sólo de los ángulos, por lo que ambos deben ser iguales a una constante que denotamos $\ell(\ell+1)$:

\textbf{Ecuación angular:}
$$\frac{1}{\sin\theta}\frac{\partial}{\partial\theta}\left(\sin\theta\frac{\partial Y}{\partial\theta}\right) + \frac{1}{\sin^2\theta}\frac{\partial^2 Y}{\partial\varphi^2} + \ell(\ell+1)Y = 0.$$

\textbf{Ecuación radial:}
$$\frac{d}{dr}\left(r^2\frac{dR}{dr}\right) + [k^2r^2 - \ell(\ell+1)]R = 0.$$


La ecuación angular es la ecuación de los armónicos esféricos. Separamos $Y(\theta,\varphi) = \Theta(\theta)\Phi(\varphi)$:

Para $\Phi(\varphi)$ obtenemos:
$$\frac{d^2\Phi}{d\varphi^2} + m^2\Phi = 0 \quad \Rightarrow \quad \Phi(\varphi) = e^{im\varphi},$$
donde $m$ debe ser entero para garantizar la periodicidad $\Phi(\varphi+2\pi) = \Phi(\varphi)$.

Para $\Theta(\theta)$ obtenemos la ecuación asociada de Legendre:
$$\frac{1}{\sin\theta}\frac{d}{d\theta}\left(\sin\theta\frac{d\Theta}{d\theta}\right) + \left[\ell(\ell+1) - \frac{m^2}{\sin^2\theta}\right]\Theta = 0.$$

Con el cambio $x = \cos\theta$, esta se convierte en la ecuación de Legendre asociada cuyas soluciones regulares en $[-1,1]$ son los polinomios asociados de Legendre $P_\ell^m(\cos\theta)$, con $\ell$ entero no negativo y $|m| \leq \ell$.

La ecuación radial es:
$$\frac{d}{dr}\left(r^2\frac{dR}{dr}\right) + [k^2r^2 - \ell(\ell+1)]R = 0.$$

Haciendo el cambio $R(r) = u(r)/r$, obtenemos:
$$\frac{d^2u}{dr^2} + \left[k^2 - \frac{\ell(\ell+1)}{r^2}\right]u = 0.$$

Esta es la ecuación de Bessel esférica. Introduciendo la variable adimensional $x = kr$, la ecuación se escribe:
$$\frac{d^2u}{dx^2} + \left[1 - \frac{\ell(\ell+1)}{x^2}\right]u = 0.$$

Las soluciones regulares en el origen son las funciones de Bessel esféricas $j_\ell(x)$, relacionadas con las funciones de Bessel ordinarias por:
$$j_\ell(x) = \sqrt{\frac{\pi}{2x}}J_{\ell+1/2}(x).$$

Por lo tanto, la solución radial es:
$$R(r) = A j_\ell(kr),$$
donde hemos descartado la función de Neumann esférica $\eta_\ell(kr)$ porque diverge en $r=0$.

La condición de que $\psi$ sea nula sobre las paredes ($r=R$) implica:
$$R(R) = A j_\ell(kR) = 0.$$

Para tener una solución no trivial ($A \neq 0$), debemos tener:
$$j_\ell(kR) = 0.$$

Denotemos por $\alpha_{\ell n}$ el $n$-ésimo cero de la función $j_\ell(x)$, es decir:
$$j_\ell(\alpha_{\ell n}) = 0, \quad n = 1, 2, 3, \ldots$$

Entonces:
$$kR = \alpha_{\ell n} \quad \Rightarrow \quad k_{\ell n} = \frac{\alpha_{\ell n}}{R}.$$

Recordando que $k^2 = 2mE/\hbar^2$, obtenemos los niveles de energía cuantizados:
$$E_{\ell n} = \frac{\hbar^2 k_{\ell n}^2}{2m} = \frac{\hbar^2}{2m}\left(\frac{\alpha_{\ell n}}{R}\right)^2 = \frac{\alpha_{\ell n}^2\hbar^2}{2mR^2}.$$

Para cada par de números cuánticos $(\ell, n)$, tenemos $(2\ell+1)$ estados degenerados correspondientes a los diferentes valores de $m = -\ell, -\ell+1, \ldots, \ell$.

La solución general es una combinación lineal de todas las soluciones permitidas:
$$\psi(r,\theta,\varphi) = \sum_{\ell=0}^{\infty}\sum_{m=-\ell}^{\ell} A_{\ell m} j_\ell(k_{\ell n}r) Y_{\ell}^m(\theta,\varphi).$$

Sustituyendo $k_{\ell n} = \alpha_{\ell n}/R$:
$$\psi(r,\theta,\varphi) = \sum_{\ell=0}^{\infty}\sum_{m=-\ell}^{\ell} A_{\ell m} j_\ell\left(\frac{\alpha_{\ell n}r}{R}\right) Y_{\ell}^m(\theta,\varphi).$$

\end{document}
